\chapter{Get, Configure and Compile CFS++}
This section provides a description for new  CFS++ developer  to  get the
necessary environment running for doing  CFS++ development. We
start with the  required software packages from the  operating system software
distribution services  and give the  required commands to install  them. After
that we describe how to get the CFS++ sources from the development server.  In
the next  section we  describe how the  source tree  has to be  configured for
building binaries.  We also mention how  to generate projects  for Eclipse. An
important point  while doing  FEM development is  pre- and  postprocessing. We
therefore also  describe how to obtain  the pre- and  postprocessors which are
supported by CFS++.  In the last section we give hints  on what other chapters
and sections are a good further reading for users with different interests.

In order  to get started with  CFS++ development it  is a good idea,  to learn
about  the   usage  of   CFS++  first  before   going  into  the   details  of
development. The  document which you  should therefore consider first,  is the
CFS++ user's tutorial \cite{Simkovics2010}. If you are comfortable using CFS++
the next  question to ask is,  wether you want to  setup up a  new machine for
doing  development, or  if  you are  going  to use  one  of the  preconfigured
machines in Erlangen or Klagenfurt.  In  the former case you should start with
sections  \ref{sec:list_of_ext_packages}  and \ref{sec:obtaining_prepostprocs}
before following the  steps outlined in the remaining  sections of the current
chapter.    In  Sec.    \ref{sec:first_time_implementing}  a   somewhat  rough
description is given about the  structure of CFS++ by explaining the necessary
code parts for  a pseudo PDE.  To get hands-on experience  with the code, Sec.
\ref{sec:code_in_action}  describes how  to examine  the static  and transient
behavior of the code by recording callgraphs and by using a debugger.

\section{Setting up a Machine for CFS++ Development}
\label{sec:list_of_ext_packages}

\subsection{UNIX-based Operating Systems}

The list of  required packages depends on the  operating system and especially
for Linux based OSes on the distribution.  We provide a bootstrap shell script
for installing the required packages  from the command line.  The script works
for some of the  OSes most often in use by the  CFS++ developers. It is called
{\tt  bootstrap\_devel\_machine.txt}  and  can  be  obtained  from  the  CFS++
development server.   The shell  script needs to  be executed with  super user
privileges (consult  your computer's administrator!).   The following sequence
of commands is required for the setup process:

\input{pygmentized/qs_bootstrap.tex}
\attachfile[icon=PushPin,description={Shell   script   for   bootstrapping   a
development machine.},mimetype=text/plain]{attachments/qs_bootstrap.txt}

If  you use a  different distribution  you should  be able  to figure  out the
required  packages  by  examining  {\tt  bootstrap\_devel\_machine.txt}  quite
easily.

The  file  {\tt  bootstrap\_devel\_machine.txt}   contained  in  this  PDF  is
automatically  generated by  concatenating  {\tt share/scripts/distro.sh}  and
{\tt  share/scripts/bootstrap\_devel\_machine.sh} from  the CFS++  source code
directory.   If  you want  to  implement  support for  a  new  OS, you  should
therefore follow the examples given there  and commit your changes back to the
repository.

\subsubsection{Linux}

The  description given in  the section  above applies  to all  supported Linux
distributions. These include:

\begin{description}

\item[openSUSE and SLE] All openSUSE, SLES and SLED versions from  the 11.x series should work out of
the box.   If some  package is not  available from the  standard repositories,
there   is  a   good  chance   to  find   it  at   openSUSE   software  search
\url{http://software.opensuse.org/search},      openSUSE     build     service
\url{https://build.opensuse.org} or one of the other package repositories like
PackMan (\url{http://packman.links2linux.org}).

\noindent Checking which package provides an already installed file:
\begin{verbatim}
rpm -q --whatprovides /usr/bin/g++
\end{verbatim}

\item[Debian Based]  Ubuntu versions  from 10.04  LTS  are known  to work.   Earlier
versions  may also  work. Since  Ubuntu is  derived from  Debian,  also recent
Debian releases should work. Support for Linux Mint, which is in turn a
spinoff from Ubuntu, has also been incorporated.

\noindent Checking which package provides an already installed file:
\begin{verbatim}
dpkg -S /usr/bin/g++
\end{verbatim}

Ubuntu is a little  special among other Linux distros since it  uses Dash as a
shell instead of  Bash. Therefore some shell scripts written  for Bash may not
work on  Ubuntu. Go  to \url{https://wiki.ubuntu.com/DashAsBinSh} for  tips on
how  to  make  your  scripts  POSIX  compliant. We  also  provide  the  script
\texttt{share/scripts/checkbashisms} which  gives warnings about bash-specific
syntax. On  Ubuntu, there exists  also the \texttt{devscripts}  package, which
contains a lot of other tools for script development.

\item[Fedora] Fedora support has been tested on release 13.

\noindent Checking which package provides an already installed file:
\begin{verbatim}
rpm -q --whatprovides /usr/bin/g++
\end{verbatim}

\item[CentOS] Support has  been added for CentOS 5.x  and 6.x. Therefore, also
other RHEL clones (e.g. Scientific Linux) should work.
\end{description}%

\subsubsection{Mac OS X}

The following list includes a number of required packages on Mac OS X. 
\begin{itemize}
\item XCode - \url{http://developer.apple.com/technologies/tools/xcode.html}
\item MacPorts - \url{http://www.macports.org}
\item CMake - \url{http://www.cmake.org}
\item g++, tcl,  etc. from MacPorts 
%\attachfile[icon=PushPin,description={List
%of                   packages                  installed                  from
%MacPorts.},mimetype=text/plain]{attachments/mac_ports_installed_pckgs.txt}
%cf. \url{http://pc.freeshell.org/comp/macbook}
\item For Fortran go to \url{http://www.macresearch.org/gfortran-leopard}. Also
visit \url{http://hpc.sourceforge.net/}.
\item Building a cross compiler for Mac OS X on Linux: \url{http://devs.openttd.org/~truebrain/compile-farm/apple-darwin9.txt}
\end{itemize} 

The  steps described  in Sec.   \ref{sec:list_of_ext_packages} also  guide the
user through  the setup  process of a  Mac OS  X development machine.   It was
tested on Snow  Leopard (10.6) but should also work on  Leopard (10.5) and the
 Lion (10.7) release.

\subsection{Windows (XP, Vista, 7)}

We support building of CFS++ for Windows in two different flavors:

\begin{description}
\item[Cross-compiling from Linux] At the moment we support two MinGW toolchains on CentOS 6 (gcc 4.5.3)
and Ubuntu 12.04 (gcc 4.6.3, cf. Sec. \ref{sec:mingw_linux_crosscompile}).
\item [Compiling on Windows] by making use of the MinGW toolchain (gcc 4.7.1) and some additional GNU utilities. 
\end{description}%
%
Unlike on UNIX systems, there is no automated setup script for these tools. Therefore, we just describe
here, where to obtain the tools and how to set up a build environment on Windows.

\begin{description}
\item[GNU Compilers] There exist several distributions for the MinGW toolchain. The one we used for testing was the 
\emph{TDM-GCC}\footnote{TDM-GCC toolchain: \url{http://tdm-gcc.tdragon.net}} toolchain. 
In their download section one can either download, single packages or a bundle installer for either 32-bit or 64/32-bit development.
We chose the tdm64-gcc bundle installer. This package provides a standard Windows installer with all required packages
for C, C++ and Fortran compilers. Another toolchain is available at \url{http://www.equation.com}.

\item[GNU Tools] We also need the usual GNU tools, like gmake, for building software. There are also two options available for obtaining these.
There exists an older project \emph{GNU utilities for Win32}\footnote{GNU utilities for Win32: \url{http://unxutils.sourceforge.net}} or
the \emph{GnuWin32}\footnote{GnuWin32 project: \url{http://gnuwin32.sourceforge.net}} project, which provides more up-to-date versions of the utilities.
After following their corresponding download and installation instructions, one should end up with the installed trees. One has to add the corresponding
binary dirs ({\tt GNUWIN32\_INSTALL\_DIR/bin} for GnuWin32 or {\tt UNXUTILS\_INSTALL\_DIR/usr/local/wbin} for UnxUtils) to the system path

\item[CMake] for Windows can be obtained from download section of the official homepage\footnote{CMake  download page: \url{http://www.cmake.org/cmake/resources/software.html}}.

\item[Subversion] The best-known Subversion distribution on Windows seems to be TortoiseSVN\footnote{Subversion for Windows: \url{http://tortoisesvn.tigris.org}}. A Java-based platform independent alternative which also provides command line tools is SVNKit\footnote{SVNKit: \url{http://svnkit.com}}.

\item[Git] If the CFS++ working copy is Git-based, then one should get Git for Windows\footnote{Git for Windows: \url{http://msysgit.github.com}}. There seems to be also a TortoiseGit project\footnote{TortoiseGit project: \url{http://code.google.com/p/tortoisegit}}.

\item[Intel MKL] If the optimized solvers from MKL need to be used, a Windows distribution (e.g. as provided by the latest Intel Compiler installations) is required. Unlike on UNIX-based systems, no advanced system for automatic detection of the required
link parameters has been implemented so far. The corresponding settings can be found and hand-tuned in {\tt cmake\_modules/FindIntelMKL.cmake}, however. The only version of MKL known to work on Win64 with MinGW at the moment is 10.0.5.025.
\end{description}%
%
The binary directories of the aforementioned distributions need to be added to the system PATH, so 
that the utilities can be found afterwards by our build system. The configuration of CFS++ on Windows is then pretty similar as compared to UNIX. The only difference is, that one should specify the additional parameter {\tt -G 'MinGW Makefiles'} on the command line of CMake, since Microsoft NMake files or Visual Studio projects would be generated otherwise.

At the moment, we just support the absolute core components of CFSDEPS and CFS++ on Windows. Items like HDFView or ParaView are still to be ported. 

There is also only very limited support for the Microsoft Visual Studio compilers, since we require also a Fortran compiler. Since a Fortran compiler is not part 
of the Microsoft tools, either the Intel Fortran compiler would be needed or the GNU Fortran compiler could be used.
 This would however, require some changes to our build project\footnote{Fortran for C/C++ developers made easier with CMake: 
\url{http://www.kitware.com/blog/home/post/231}}. These changes may be worth the effort in the future, since Visual Studio Express is 
an excellent environment for debugging and is also available for free since the 2008 release.


\todo[inline] {
Pre-installed   VBoxes  containg   all  required   packages  are   located  at
(Vienna)}

\noindent {\tt qnap:/share//AM\_Data/Software/NACS\_Windows\_Development\_VBOXes/[WINXP*\&nacs\_dev*\&local*]}.

%Refer   to   the   howto  \attachfile[icon=PushPin,description={NACS   Windows
%development           howto           on           dedicated           Virtual
%Boxes.},mimetype=text/plain]{attachments/nacs_win_build_howto.txt} for getting
%these up and running.

\newpage
\section{Obtaining the Sources}

Once the software  needed for CFS++ development from  the distribution has
been installed,  the next  step is  to checkout the  source code  from the
Subversion  \cite{svn} repository. All  necessary steps  are given  in the
shell  script at  the end  of this  section. We  first create  a directory
structure for our  development resources. Next we obtain  the sources from
the server. Note,  that upon the first connection you  will have to accept
the server's  certificate permanently (p)  and decide, wether you  wish to
save your password locally.

\input{pygmentized/qs_obtain_sources.tex}
\attachfile[icon=PushPin,description={Shell script for obtaining CFS++ sources.},mimetype=text/plain]{attachments/qs_obtain_sources.txt}

Further information  about Subversion can be found  in Sec. \ref{sec:svn}.

A nice alternative, especially for working on a laptop, is the distributed
versioning system Git.  A basic howto for obtaining a  Git working copy of
the CFS++ sources, some basic commands and pointers to additional material
are given in Sec. \ref{sec:git}.

\section{Configuring the Build Tree and Building Binaries}

We use  the CMake build system  (cf.  Sec.  \ref{sec:cmake})  to configure our
source tree for building on a  wide range of platforms.  Unlike the well-known
configure/make/make install methodology from Unix-like systems, CMake supports
also other  platforms like Windows.   It generates build descriptions  for the
native  development environments (e.g.   Makefiles on  Unix and  Visual Studio
projects  on Windows)  from very  simple input  files in  CMake  syntax. Using
CMake,  it   is  also   possible  to  do   out-of-source  builds   (cf.   Fig.
\ref{fig:understand_src_bin_dir}), which means, that from a single source tree
multiple   build   trees   may   generated   (e.g.  for   debug   or   release
executables). This way,  CMake keeps the different builds  separate and avoids
cluttering the  source directory with a  huge number of  generated files.  The
latter behavior can  make the output of the  Subversion status commands almost
useless.  By  doing dependency  analysis  CMake  can  also generate  Makefiles
suitable for parallel builds.

\begin{figure}[htb] 
  \centering
  \begin{overpic}[width=0.8\textwidth]{pics/understand_src_bin_dir}
  \end{overpic}
  \caption{Two  build directories  with  different settings  generated from  a
single source directory.}
  \label{fig:understand_src_bin_dir}
\end{figure}%

In the following script we create  two empty build directories. In one of them
we will  build a debug executable,  whereas in the  other one we will  build a
release  executable.   Before configuring  we  set  our  desired compilers  by
creating  new  environment variables.   After  that  we  change to  the  build
directories and configure them.  For this  purpose we also have to specify the
location of CFSDEPS to CMake and where the dowloaded sources of libraries will
be  stored  (CFS\_DEPS\_CACHE\_DIR, cf.   Sec.   \ref{sec:cfsdeps}).  We  also
specify, that we want to use standard  Unix make files for building. As a last
step  we  have to  build  the  libraries, that  CFS++  depends  upon, and  the
executables.  This step  may take a considerable amount of  time for the first
time, but  will not  be necessary anymore,  if the  CFSDEPS do not  change any
more. Please  note, that you should always  specify the build of  cfsdeps in a
separate step before the build of CFS++, since race conditions during parallel
builds could occur at the moment.   If your home directory is a network share,
please  take  Sec.   \ref{sec:being_careful_with_shares} into  account  before
executing the following commands.

\input{pygmentized/qs_configure_build}

\subsection{Being Careful with Network Shares}
\label{sec:being_careful_with_shares}
\textcolor{red}{IMPORTANT  NOTE:  Avoid  putting  your  build  directories  on
network shares! User home directories tend  to be such shares.  CMake and make
require certain  file access mechanisms, that  may not be  available on shared
folders. Moreover,  the build will take much  longer and it may  not even make
sense to share executables  among different machine architectures or operating
system versions!}

You can check  whether your home directory is mounted from  a network share by
typing
\begin{verbatim}
mount | grep $HOME
\end{verbatim}
into a shell prompt. If you get a response like this
\begin{verbatim}
//server/user_homes/georg/ on /home/users/georg type cifs (rw)
\end{verbatim}
you should not put your build directory inside your home directory.

It is however perfectly sane to  put your source directory on a network share.
Such practice is even advisable, since changes need to be made only once for a
set   of   sources   shared   by   multiple   computers.    Refer   to   Chap.
\ref{chap:mounting_sshfs}.  for sharing  the source directories with computers
which only have SSH access to the machine containing the sources.

\section{Generating a Project for Eclipse}
\label{sec:eclipse_setup_project}

The only  difference from the  procedure for generating Makfiles  described in
the section  above, is  that we  need to specify  a different  generator while
configuring the build tree with CMake. We do this by saying:

\begin{verbatim}
cmake -DDEBUG:BOOL=ON \ 
      -DCFS_DEPS_CACHE_DIR:PATH=$DEVDIR/CFSDEPSCACHE \ 
      -G "Eclipse CDT4 - Unix Makefiles" 
      $DEVDIR/CFS_TRUNK
\end{verbatim}

Now we have generated the two  files {\tt .cproject} and {\tt .project} in the
directory {\tt  \$DEVDIR\-/build\-/debug} in addition to the  Unix make files.
In Eclipse (with the CDT extension) we have to import our new project by doing
the   following  steps:  File   $\rightarrow$  Import   $\rightarrow$  General
$\rightarrow$  Existing  Project  into  Workspace  $\rightarrow$  Select  Root
Directory {\tt \$DEVDIR\-/build\-/debug}.

In order for the Eclipse code indexer  to find all headers it may be necessary
to add some  system and MKL include paths and  preprocessor defines.  This can
be  done under  Project $\rightarrow$  Properties $\rightarrow$  C/C++ Include
Paths and Symbols (cf.  Fig. \ref{fig:eclipse_cpp_inc_paths} for an example on
Mac OS X).

\begin{figure}[htb] 
  \centering 
  \subfloat[Additional include paths.]{
    \label{fig:eclipse_cpp_inc_paths}
    \begin{overpic}[height=5cm]{pics/eclipse_cpp_inc_paths}
    \end{overpic}
    }
  \qquad
  \subfloat[Importing a coding style.]{
    \label{fig:eclipse_import_bsd_style}
    \begin{overpic}[height=5cm]{pics/eclipse_bsd_style}
    \end{overpic}
  } 
  \caption{Specifying additional include directories and coding style settings.}
  \label{fig:eclipse_code_style_inc}
\end{figure}

In order to  maintain a common coding style  among Eclipse developers everyone
should     import    the     style    informations     attached     to    this
PDF.\attachfile[icon=PushPin,description={Eclipse                        coding
style.},mimetype=text/plain]{attachments/eclipse-code-style.txt}. The settings
can be  imported into  Eclipse by going  to Project  $\rightarrow$ Preferences
$\rightarrow$ C/C++  General $\rightarrow$ Code  Style $\rightarrow$ Configure
Workspace             Settings...              $\rightarrow$            Import
(cf. Fig. \ref{fig:eclipse_import_bsd_style}).

Also the  settings for  character encoding and  further editor  setting should
remain consistent among the developers.  The necessary settings in Eclipse can
be made  by accessing  Window $\rightarrow$ Preferences  $\rightarrow$ General
$\rightarrow$ Editors and Workspace (c.f Fig. \ref{fig:eclipse_cfs_settings}).
\begin{figure}[htb] 
  \centering 
  \subfloat[Editor settings.]{
    \label{fig:eclipse_editor_settings}
    \begin{overpic}[height=5cm]{pics/eclipse_editor_settings}
%      \put(10,20){\textcolor{red}{\bf Test}}
    \end{overpic}
    }
  \qquad
  \subfloat[Workspace settings.]{
    \label{fig:eclipse_workspace_settings}
    \begin{overpic}[height=5cm]{pics/eclipse_workspace_settings}
    \end{overpic}
  } 
  \caption{Settings for CFS++ in Eclipse.}
  \label{fig:eclipse_cfs_settings}
\end{figure}

For   convenience  the   corresponding  settings   files  are   attached  here
\attachfile[icon=PushPin,description={Eclipse                            editor
settings.},mimetype=text/plain]{attachments/org.eclipse.ui.editors.txt}
\attachfile[icon=PushPin,description={Eclipse                         workspace
settings.},mimetype=text/plain]{attachments/org.eclipse.core.resources.txt}. The
have               to                copied               to               the
workspace/.metadata/.plugins/org.eclipse.core.runtime/.settings  directory and
must be renamed to \*.prefs.


\section{Pre-/Postprocessors}
\label{sec:obtaining_prepostprocs}

Since its  creation CFS++  has come to  support a  growing number of  pre- and
postprocessors.  The three  most  important  being GiD, Gmsh  and  ANSYS Classic.  The
modelling process with these tools is documented in \cite{Simkovics2010}. 

For the broader  picture of the software landscape around  CFS++ the reader is
referred    to     the    Advanced    Pre-     and    Postprocessing    Manual
\cite{CFSAdvPrePostManual}, which  features information on  common and special
pre-  and  postprocessing tasks  in  recipe-like style  for  a  wide range  of
commercial and free pre- and postprocessors.

\section{The CFS++  Development Server}

All necessary resources (e.g. Subversion  server, Trac, CDash, etc.) for doing
CFS++ development have been installed  on common a dedicated server machine at
Erlangen.  Also a web server is running there, which serves as a portal to the
different       services.        It        can       be       reached       at
\url{https://lse17.e-technik.uni-erlangen.de:2001}. In this section we briefly
describe the services offered by the server.

\begin{figure}[htb] 
  \centering
  \begin{overpic}[width=0.8\textwidth]{pics/cfs_devel_server}
  \end{overpic}
  \caption{Home page of CFS++ development server with links to different services.}
  \label{fig:cfs_devel_server}
\end{figure}%

Figure \ref{fig:cfs_devel_server} shows the home page of the server along with
the pages for the Trac system  (cf. Sec. \ref{sec:trac}) and the CDash nightly
build  dashboard  (cf.  Sec.   \ref{sec:cdash}).   Furthermore the  Subversion
repository  can be browsed  online using  the SVN  link.  Nightly  built CFS++
manuals, Doxygen  documentation and  other documents related  to CFS++  can be
accessed via the  Documentation link. Under File Management  the developer can
find  nightly  built and  other  pre-built  CFS++  executables. Also  GMV  and
ParaView  executables  for different  platforms  can  be  found there.   Users
registered as administrators  may change the server settings  by accessing the
Administration link.


